\chapter{内部热阻不可辨识性、对偶探针与安全屏障}
\label{ch:dual}

上一章把连续方程收缩成十六条控制体和一个两节点降阶，硬件上也就能算出冷点该怎么走。但收缩之后还有一个更要紧的问题没回答：只靠腔体侧能装上的那些传感器，这一次过程里的内部热阻究竟能不能被看见。

直觉上好像可以：腔温采样能做到很密，表面温度也能测得很准，把曲线画细一些，总该能推断出内部发生了什么。问题在于平衡点上的代数结构不允许这样做。当食品与腔体都停在稳态，内部热阻、表面热容与内部热容只以除数形式出现在方程里，求导之后右端为零；换言之，稳态观测对这三个量的灵敏度严格为零，采样再密，也只是把一条已经闭合的空气温度曲线画得更密。

出路只能是离开平衡。本章围绕这个结论安排三件事。其一，证明平衡观测对内部热阻的灵敏度为零，把“为什么稳态测量不起作用”写成命题。其二，给出一条可行的激励方案：家电嵌入式计算能力有限，在线求 Fisher 矩阵不现实，改用固定时窗的两段开环激励，先蒸汽、后风机，从而把外部换热与内部导热的两列灵敏度分开。其三，处理这样做的代价：蒸汽对品质的作用不止一条，它同时经冷点潜热、表面蒸发与慢热阻三条路径影响结果，必须分开计算；而安全投影只读水位、门、剩余电流与加热器温度，品质估计不得进入断电条件。三件事合起来，才构成从“看不见”到“看得见且不越界”的完整过渡。

\section{可安装观测与高频噪声的放大}
\label{ch:est}

家用蒸烤箱上已经装得上的通道其实不少：腔温热电偶在后壁或顶部，湿度探头靠近排汽口，门缝处可以加一根穿刺针，部分机型还会在门玻璃内侧装红外。表\ref{tab:obs}把这些通道逐条列了出来，其中只有一部分真正进入品质估计器——安全侧的测量（水位、门位、剩余电流与加热器温度）根本不进品质滤波。腔温、湿度、质量和红外都能写成观测向量的分量，而工程核实际读入的只有两条：表面温度加 \(0.4^\circ\mathrm{C}\) 的确定性散列，以及平均含水加 \(0.004\) 的散列。其余通道留在表里，正是为了标出“可安装”与“已经接到跟踪器”之间的那段距离。

\begin{table}[htbp]
\centering
\caption{进入品质估计器的观测}
\label{tab:obs}
\setlength{\tabcolsep}{2.2pt}
{\footnotesize
\begin{tabularx}{\columnwidth}{llY}
\toprule
通道 & 符号 & 与隐藏状态的关系 \\
\midrule
腔体空气温度 & \(y_T\) & 靠近 \(T_\infty\)，受加热器辐射污染 \\
腔体相对湿度 & \(y_\varphi\) & \(p_v/p_{\mathrm{sat}}(T_\infty)\)，饱和时丢失分辨率 \\
食品质量 & \(y_m\) & 累计蒸发的积分，不定位冷点 \\
表面热像 & \(\vect{y}_{\mathrm{IR}}\) & 发射率未知的辐射亮度 \\
电功率 & \(y_P\) & 已知输入，写入边界 \\
排气温度 & \(y_{T,e}\) & 冷凝与漏汽的积分证据 \\
\bottomrule
\end{tabularx}}
\end{table}

装得上并不等于已经看见内部热阻，这中间的落差由每一路通道各自的缺陷决定。腔温离空气太近，容易被加热器的辐射污染；湿度在饱和段几乎没有分辨率可言；质量是累计蒸发的积分量，根本无法定位冷点；红外读到的是发射率未知的辐射亮度。穿刺针虽然能靠近内部，却会改变局部场，而且一旦没刺中冷点，读数就不代表场的下确界。工程核因此只用表面温度与平均含水这两条通道，再加一个确定性散列。探针窗的任务就清楚了：让这两条已经装得上的通道暂时离开平衡，从而把内部热阻从灵敏度为零的状态里拖出来。

观测方程 \(\vect{y}_k=h\big(\Phi(\vect{X}(t_k)),\vect{\theta}\big)+\vect{v}_k\)。红外通道
\[
L_j=\epsilon(\varepsilon_{l,s},S_s)\,\sigma T(\vect{r}_j)^4
+\big(1-\epsilon(\varepsilon_{l,s},S_s)\big)L_j^{\mathrm{refl}}.
\]
发射率随表面变干下降。系统精度应取发射率区间两端分别滤波后的分位包络。进入跟踪器的噪声由种子和步号生成，重复运行得到同一轨迹。

线性化热方程在厚度方向上，对高频边界噪声的放大随 \(\mathrm{e}^{kL}\) 增长\cite{ozisik2000}。半无限体上表面温度取成角频率 \(\omega\) 的谐波时，
\begin{equation}
\label{eq:halfspace}
T(x,t)=\mathrm{e}^{-kx}\cos(\omega t-kx),\qquad
k=\sqrt{\frac{\omega}{2\alpha}}.
\end{equation}
深度 \(L\) 处的振幅是表面振幅的 \(\mathrm{e}^{-kL}\)。\(\alpha=2.0\times 10^{-7}\,\mathrm{m}^2/\mathrm{s}\)。表\ref{tab:amp}给出正文中已经出现的时间尺度。

\begin{table}[htbp]
\centering
\caption{半无限体上单个谐波的放大倍数 \(\mathrm{e}^{kL}\)}
\label{tab:amp}
\setlength{\tabcolsep}{2.6pt}
{\footnotesize
\begin{tabular}{llrr}
\toprule
周期 & 位置 & \(L=30\,\mathrm{mm}\) 放大 & 穿透深度 \\
\midrule
8 s & 演示记录间隔 & \(1.81\times 10^{18}\) & 0.71 mm \\
40 s & 探针窗 & \(1.46\times 10^{8}\) & 1.60 mm \\
1170 s & \(30\,\mathrm{mm}\) 开环菜单 & 32.3 & 8.63 mm \\
\bottomrule
\end{tabular}}
\end{table}

以 \(8\,\mathrm{s}\) 的谐波为例，传到 \(15\,\mathrm{mm}\) 处只剩下 \(7.44\times10^{-10}\,\mathrm{K}\)，放大到这个程度已经没有任何信息可用。菜单时间又把放大压到另一个量级：薄件在自己的菜单时刻还留着约十分之一，厚件只剩千分之几。若令 \(kL=1\)，对应的周期是 \(\pi L^2/\alpha\)，\(30\,\mathrm{mm}\) 上算得 \(14137\,\mathrm{s}\)，比整条菜单还长。由此可见，加密表面像素所能提高的只是那些早就衰减掉的高频分量；要真正看见冷点，必须提供穿透时间尺度上的动态激励，或者找一个对内部热阻敏感的积分量。质量损失恰好两样都不占：它对累计蒸发敏感，却对冷点位置不敏感。

\section{稳态观测为什么不可辨识}

\begin{proposition}[平衡点不含内部参数]
\label{prop:steady}
设\eqref{eq:twonode}中 \(\dot m=0\)、\(q_{\mathrm{lat}}=0\)，且存在平衡点。则 \(T_1=T_2=T_\infty\)。\((R_{\mathrm{int}},C_1,C_2)\) 不进入该代数方程，平衡观测对三者的灵敏度为零。
\end{proposition}

\begin{proof}
第二式在平衡时为 \((T_1-T_2)/R_{\mathrm{int}}=0\)。热阻取正，故 \(T_1=T_2\)。第一式变为 \((T_\infty-T_1)/R_{\mathrm{ext}}=0\)，故 \(T_1=T_\infty\)。三个参数只作为除数出现，不出现在解的表达式中。对任一参数求导，平衡方程的右端仍为零。
\end{proof}

由此得到一个直接推论：恒温段里腔温采得再密，也不会更新内部热阻，信息动作必须离开平衡。对 \(R\in\{1,3,6\}\) 从同一初值积分，平衡温差的极差恰好为 \(0\)。这正是“到时按菜单停机”失败的数学形式——菜单时钟闭合的是时间，而不是这一次过程里被看见的 \(R_{\mathrm{int}}\)。

上面的证明给出平衡点的结论，但还需要说明一件事：为什么"离开平衡"就能恢复可辨识性。这要看灵敏度矩阵在平衡与准稳态两处的秩。

\begin{proposition}[平衡点的灵敏度为零，偏离后立即恢复满秩]
\label{prop:rank}
设观测为表面温度 \(T_1\)，待估参数为 \((R_{\mathrm{int}},R_{\mathrm{ext}})\)。

(i) 在平衡点 \(T_1=T_2=T_\infty\) 处，灵敏度矩阵
\(S=\partial T_1/\partial(R_{\mathrm{int}},R_{\mathrm{ext}})\) 的两列同时为零，秩为 \(0\)。

(ii) 一旦离开平衡，只要净热流 \(q\neq0\)，\(S\) 即满秩。
\end{proposition}

\begin{proof}
(i) 由命题\ref{prop:steady}，平衡时 \(q=0\)，\(T_1\) 对两个热阻的偏导都含因子 \(q\)，故同时为零。

(ii) 准稳态下 \(T_1=T_\infty-qR_{\mathrm{ext}}\)，且 \(q=(T_\infty-T_2)/(R_{\mathrm{int}}+R_{\mathrm{ext}})\)。
对两个参数求偏导得
\[
\frac{\partial T_1}{\partial R_{\mathrm{int}}}
=-R_{\mathrm{ext}}\frac{\partial q}{\partial R_{\mathrm{int}}},\qquad
\frac{\partial T_1}{\partial R_{\mathrm{ext}}}
=-q-R_{\mathrm{ext}}\frac{\partial q}{\partial R_{\mathrm{ext}}}.
\]
当 \(q\neq0\) 时第二式含独立项 \(-q\)，两列不成比例，故满秩。
\end{proof}

这个结论有一个可以直接检验的推论：可辨识性是**开关式**的，在平衡点恰好为零，离开平衡就有，不需要等到某个"足够大"的激励幅度。数值上，把 \(R_{\mathrm{int}}\) 取 \(1,3,6\) 在准稳态分支上求偏导，堆叠后奇异值为 \(117.96\) 与 \(40.68\)，秩为 \(2\)；而在平衡点同一矩阵恒为零，秩为 \(0\)。两者之间没有过渡带，原因就是热流 \(q\) 在平衡点处为零。命题\ref{prop:steady} 说的是"看不见"，命题\ref{prop:rank} 补上了"什么时候开始看得见"，两段探针的设计依据就在后者：它只要让 \(q\) 离开零，并且在两个参数上给出不同的偏导路径。



图\ref{fig:inobserv}把这条命题画成一段可以看的时序，三个内部热阻从同一个平衡点出发，在前 \(30\,\mathrm{min}\) 里保持在 \(60^\circ\mathrm{C}\) 不动——平衡段冷点的三档差小于 \(0.01\,\mathrm{K}\)，任何采样密度都分不开它们；边界在 \(30\,\mathrm{min}\) 跳到 \(150^\circ\mathrm{C}\) 之后，三条曲线立刻按 \(R\) 排序分开，\(R\) 越大冷点越靠下。命题说的是代数事实：平衡解 \(T_1=T_2=T_\infty\) 里三个参数只作为除数出现，求导后右端仍为零。图说的是同一件事的时序版本：**信息不在平衡点上，而在离开平衡的方向上**。

\begin{figure*}[t]
\centering
\spanimg{X_inobserv.png}
\caption{平衡段三条不同热阻的轨迹完全重合（左：表面实线、冷点虚线），边界激励后按 \(R\) 排序分开（右）。这是命题\ref{prop:steady} 的时序版本}
\label{fig:inobserv}
\end{figure*}

灵敏度 \(\vect{\sigma}=\partial(T_1,T_2)/\partial R_{\mathrm{int}}\) 满足
\begin{equation}
\label{eq:sigma}
\dot{\vect{\sigma}}=\vect{A}\vect{\sigma}+\frac{\partial\vect{A}}{\partial\theta}\begin{pmatrix}T_1\\ T_2\end{pmatrix}.
\end{equation}
平衡且 \(T_1=T_2\) 时右端为零，\(\vect{\sigma}=0\) 是平衡解，Fisher 矩阵里这一列贡献为零。离开平衡、温差不为零时，右端不再是零。

参数灵敏度满足
\begin{equation}
\label{eq:sens}
\frac{\mathrm{d}}{\mathrm{d}t}\frac{\partial \vect{s}}{\partial\theta_i}
=\vect{A}(\vect{\theta})\frac{\partial \vect{s}}{\partial\theta_i}
+\frac{\partial\vect{A}}{\partial\theta_i}\vect{s}
+\frac{\partial\vect{B}}{\partial\theta_i}\vect{u}.
\end{equation}
Fisher 信息
\begin{equation}
\label{eq:fisher}
\vect{\mathcal{I}}(\vect{\theta};\vect{u}_{0:K})
=\sum_{k=0}^{K}
\Big(\frac{\partial h}{\partial \vect{s}}\frac{\partial \vect{s}_k}{\partial\vect{\theta}}\Big)^{\!\top}
\vect{R}_k^{-1}
\Big(\frac{\partial h}{\partial \vect{s}}\frac{\partial \vect{s}_k}{\partial\vect{\theta}}\Big).
\end{equation}
探针设计在规格上最大化 \(\log\det\vect{\mathcal{I}}\) 的增量。工程核并不在每个时间步组装这个矩阵，而把同一几何压成固定时窗，见第\ref{ch:ati}节。

在线一侧读表面温度和平均含水。内部时间尺度
\[
\tau=\mathrm{clip}(\hat r,0.8,8)\,(L/0.03)^2\,280\ \mathrm{s}.
\]
估计冷点
\begin{align}
\hat T_c &\leftarrow \hat T_c+\Delta t\,(y_T-\hat T_c)/\tau,\\
\hat T_c &\leftarrow \min\big(\hat T_c,\, y_T-4\max(\hat r-1,0.2)\big).
\end{align}
第二式把估计压到表面读数之下。先验 \(\hat r=5\)，方差 \(0.64\)。\(30\,\mathrm{mm}\) 上 \(\hat r=5\) 时 \(\tau=1400\,\mathrm{s}\)；\(50\,\mathrm{mm}\) 上同一先验把 \(\tau\) 放大到约 \(3889\,\mathrm{s}\)，已经超过家用时域。先验把估计压到表面读数之下 \(4\times(5-1)=16\,\mathrm{K}\)，开局估计可以低于 \(0^\circ\mathrm{C}\)。滞后上限只用表面减一项只依赖 \(\hat r\) 的常数，表面已经远高于冷点之后，估计仍可以停在冷点之上。已实现的跟踪器给出带滞后的点估计，接口比九维决策状态窄。

\section{两阶段探针如何换取可辨识性}
\label{ch:ati}

传感器只能装在腔体一侧，这是硬件前提。一旦稳态观测压不低决策量的后验不确定度，控制器就只剩一条路：离开平衡，用短时的边界激励把内部参数从灵敏度为零的状态里拖出来。探针因此必须是一段有限时长的执行器轨迹，而不是任意一次功率变化。

\begin{definition}[热信息探针]
探针是一段有限时长 \(\tau\le\tau_{\max}\) 的执行器轨迹 \(a\in\cA\)，满足三条：设计目标是增大\eqref{eq:fisher}的 \(\log\det\)，或等价地减小决策量的后验熵；对表面褐变、质量损失和热力进度的预测损伤低于预算 \(\eta\)；结束后执行器回到探针前的名义轨迹，除非安全层已经改写该轨迹。
\end{definition}

不满足第一条的功率变化是加工动作。高斯近似下，决策状态的微分熵由协方差 \(\vect{\Sigma}_z\) 的行列式决定。对候选探针 \(a\)，
\begin{equation}
\label{eq:ati}
G(a)=\frac12\log\det\vect{\Sigma}_{z,k}
-\frac12\log\det\vect{\Sigma}_{z,k+\tau}(a).
\end{equation}
工程核用时间窗代替在线 \(\log\det\)：\([150,190)\) 秒且方差大于 \(0.25\) 时进入探针模式。前 \(20\,\mathrm{s}\) 蒸汽全开、排汽关小，后 \(20\,\mathrm{s}\) 风机升高到 \(1.3\)。顺序保留了先冷凝、后对流，时长固定为 \(40\,\mathrm{s}\)，打分用热阻是否更接近参照来代替 \(\log\det\)。

图\ref{fig:probetl}把种子 7 的执行器轨迹摊成一条时间轴，并把 \(150\)--\(190\,\mathrm{s}\) 的探针窗放大成插图看。在主轴上，探针窗只是一条四十秒宽的窄条，窗外蒸汽回到基线 \(0.15\)、风机停在 \(0.4\)、加热请求由冷点估计给出；窗内则先让蒸汽跳到 \(1.0\)，随后把风机升到 \(1.3\)。这张时间表本身就是“可分离”的实现方式——任何一个时刻，都能从图上读出当下的指令究竟出自信息动作还是品质动作。

两段为什么要换方向，可以按灵敏度看。图\ref{fig:probesens}把表面读数对内部热阻的归一化灵敏度画成三段：稳态段为零，冷凝段升到 \(0.62\)，对流段到 \(1.0\)。单靠冷凝段，\(\partial T_s/\partial R_{\mathrm{int}}\) 与外部换热的分量仍然共线；加上对流段之后，两列才线性无关，这正是命题\ref{prop:pulse} 要的第二个方向。图里的三个数不是拟合出来的系数，而是“第一段只动相变、第二段只动 \(R_{\mathrm{ext}}\)”这条设计在灵敏度上的投影。

\begin{figure}[htbp]
\centering
\colimg{M_probe_timeline.png}
\caption{种子 7 的四条执行器轨迹与探针窗（插图放大 \(150\)--\(190\,\mathrm{s}\)）。窗外蒸汽回到 \(0.15\)、风机停在 \(0.4\)；窗内先蒸汽后风机}
\label{fig:probetl}
\end{figure}

\begin{figure}[htbp]
\centering
\colimg{X_probe_sens.png}
\caption{表面读数对内部热阻的归一化灵敏度，按段给出。稳态为零；冷凝段与对流段各提供一列，拼起来才满秩}
\label{fig:probesens}
\end{figure}

窗结束时，若新的热阻比旧值更接近参照，方差乘 \(0.35\)。参照
\[
r_{\mathrm{proxy}}=\mathrm{clip}(\Delta T/8+1.2+2.5 S,\,0.8,\,20)
\]
用了表面与冷点的真值温差，实施到机器上时改成可安装通道。方差没有下降，停止时刻就留在动作之前。厚度超过 \(40\,\mathrm{mm}\) 且方差仍高于 \(0.45\) 时，停机判据要求方差先下降，点估计单独越过阈值仍保持加热。

跟踪器所估计的 \(r\) 是一个滞后倍数，即 \(\tau=r(L/0.03)^2\cdot 280\,\mathrm{s}\) 里的那个 \(r\)，而不是膜热阻与内热阻之比。它的真值可以从场上反解：在表面取一阶滞后的准稳态关系 \(\Delta T\approx\tau\,\mathrm{d}T_s/\mathrm{d}t\)，对三条厚度做最小二乘单指数拟合，对偶策略给出 \(15\,\mathrm{mm}\) 的 \(8.19\)、\(30\,\mathrm{mm}\) 的 \(8.86\)、\(50\,\mathrm{mm}\) 的 \(10.70\)，开环在 \(50\,\mathrm{mm}\) 上最高到 \(13.79\)。这些数与种子无关：同一厚度上八个种子给出同一个值，因为滞后由几何与物性决定。相比之下，参照式 \(\Delta T/8+1.2+2.5S\) 在整个定义域内只走到 \(1.24\)--\(19.35\)。若把 \(\mathrm{clip}\) 的上限放在 \(8.0\)，\(30\,\mathrm{mm}\) 对偶会有 \(48.8\%\)、\(50\,\mathrm{mm}\) 对偶会有 \(97.2\%\) 的时间贴在这个上限上；接受判据 \(|r_{\mathrm{new}}-r_{\mathrm{proxy}}|<|\hat r-r_{\mathrm{proxy}}|\) 在那些时刻实际上是在与一个常数比较。把上限放宽到 \(20\) 之后，饱和率归零，而主对照的 \(19440\) 行逐行不变（最大差 \(5\times10^{-5}\)），四个策略的中位差同样在 \(5\times10^{-5}\) 以内。这条检验记录在 \texttt{scripts/robust\_proxy.py} 与 \texttt{results/logs/proxy\_robustness.json}：参照的上限确实曾经过紧，而报告的全部结论对这处修正不敏感。

\begin{proposition}[两段探针改变的是条件数，而非秩]
\label{prop:pulse}
对\eqref{eq:twonode}在准稳态下取观测量为表面温度 \(T_1\)，参数为
\((R_{\mathrm{ext}},R_{\mathrm{int}})\)。则灵敏度矩阵在任意单一工作点上已是满秩，
两段探针的作用是把工作点沿一条使两个参数方向夹角增大的路径移动，从而降低该矩阵的条件数。
\end{proposition}

\begin{proof}
准稳态下 \(T_1=T_\infty-qR_{\mathrm{ext}}\)、\(q=(T_\infty-T_2)/(R_{\mathrm{int}}+R_{\mathrm{ext}})\)，
两个偏导为
\[
\frac{\partial T_1}{\partial R_{\mathrm{int}}}=-R_{\mathrm{ext}}\frac{\partial q}{\partial R_{\mathrm{int}}},\qquad
\frac{\partial T_1}{\partial R_{\mathrm{ext}}}=-q-R_{\mathrm{ext}}\frac{\partial q}{\partial R_{\mathrm{ext}}}.
\]
当 \(q\neq0\) 时第二式含独立项 \(-q\)，故两列不成比例，矩阵满秩。两列之比
\(\partial T_1/\partial R_{\mathrm{int}}:\partial T_1/\partial R_{\mathrm{ext}}\)
随 \(R_{\mathrm{ext}}\) 变化：在 \(R_{\mathrm{ext}}=1.0\) 处为 \(-1:3\)，
在 \(R_{\mathrm{ext}}=2.0\) 处为 \(-2:3\)。由于该比值不是常数，扫过工作点可使两列张角增大，
条件数因而下降。
\end{proof}

数值上，沿 \(R_{\mathrm{ext}}\in[0.1,2.0]\) 取十二个工作点堆叠灵敏度，
奇异值为 \(65.5\) 与 \(11.3\)，条件数 \(5.82\)，秩恒为 \(2\)。
这组数字记录在 \texttt{scripts/robust\_proxy.py} 的同源核对里。

命题\ref{prop:pulse} 修正了本节早先的一处表述：单段激励并不会让秩从 \(1\) 变成 \(2\)。
真正发生的事有两件。其一，离开平衡使 \(q\neq0\)，这正是命题\ref{prop:rank} 的内容，
秩从 \(0\) 变成 \(2\)。其二，两段探针把工作点推到两个不同的 \(R_{\mathrm{ext}}\) 上，
使灵敏度列的张角增大、条件数下降，估计对噪声更不敏感。第二件事是量的改善，
不是有无的差别，这一点必须在正文里说准，否则读者会以为一段激励完全无效。

两节点灵敏度分析给出的数值与这一读法一致：加上强对流段之后，内部热阻滞后从
单蒸汽激励下的 \(38.91\,\mathrm{s}\) 走到 \(40.85\,\mathrm{s}\)，相对倍数 \(1.05\)，
增量来自第二段对流切换，而不是潜热变号。第一段改相变符号以离开平衡，
第二段改 \(R_{\mathrm{ext}}\) 以改善条件，两段各自承担一件不同的事。

\section{对偶决策与四种停机几何}

加热请求
\begin{equation}
\label{eq:ureq}
u_{\mathrm{req}}=\mathrm{clip}\big(0.35+(75+\mu-\hat T_c)/30,\,0.15,\,1\big),
\end{equation}
裕量 \(\mu=6+10\sqrt{\mathrm{Var}}\,(L/0.03)\)。式子只看见估计温度、方差和厚度。停机要求 \(\hat T_c>75+\mu\) 且时刻已经超过 \(240\,\mathrm{s}\)。

四条策略共用同一积分器和同一条设计线 \(\alpha_m=0.85\)，差别写在停机几何上。
\begin{enumerate}
\item \textbf{开环。} 按厚度给菜单时间，中途不改。\(15\,\mathrm{mm}\)、\(30\,\mathrm{mm}\)、\(50\,\mathrm{mm}\) 分别是 \(900\,\mathrm{s}\)、\(1170\,\mathrm{s}\)、\(1530\,\mathrm{s}\)。
\item \textbf{腔温比例。} \(u=\mathrm{clip}(0.04(160-y_c)+0.6,0,1)\)。只跟踪带噪腔温，不停机，过程一直走到时域终点。
\item \textbf{无探针。} 使用\eqref{eq:ureq}，没有蒸汽窗。跟踪器看见表面温度，按点估计停机。
\item \textbf{对偶。} 在无探针加热律上叠加固定探针窗，厚件上另有方差否决。\(150\)--\(190\,\mathrm{s}\) 蒸汽和风机按时间表离开平衡，窗后若新热阻更接近参照则方差乘 \(0.35\)。
\end{enumerate}

无探针与对偶共用同一条饱和线性加热律，两者的差别只落在 \(40\,\mathrm{s}\) 窗口以及其后的热阻更新上。开环把停止时刻写成厚度的函数，腔温比例干脆把它取成时域终点 \(3300\,\mathrm{s}\)，于是四条策略在同一积分器上给出四种互不相同的停机几何：时钟、腔温、点估计，以及点估计加一段短窗。

加热请求\eqref{eq:ureq} 里的裕量随方差与厚度一起放大。方差还停在前验 \(0.64\) 时，\(30\,\mathrm{mm}\) 上 \(\mu\approx14\,\mathrm{K}\)，停机门槛大约落在 \(89^\circ\mathrm{C}\)；等方差降到 \(0.224\)，\(\mu\) 也随之退到约 \(11\,\mathrm{K}\)。厚件上同一公式会把门槛继续抬高，点估计于是更难单独宣布停机。倘若导热时间尺度在家用时域内根本走不完，停机时刻就会先被时域终点截住——这属于跟踪器一侧的几何，而不是探针拒绝工作。

信息动作与品质动作在时间轴上分开。\(150\,\mathrm{s}\) 之前四个执行量按品质加热走，加热请求由估计温度、方差和厚度决定；\(150\)--\(190\,\mathrm{s}\) 蒸汽和风机按时间表离开平衡，加热请求仍由同一条饱和律给出，只是边界被探针改写；\(190\,\mathrm{s}\) 之后蒸汽回到基线 \(0.15\)，风机回到 \(0.4\)，方差若已下降则停机门槛跟着下降。日志里要能查出这一时刻是哪一个函数发出的指令：探针窗内蒸汽列跳到 \(1.0\)，窗外蒸汽列回到 \(0.15\)，加热列始终由\eqref{eq:ureq} 给出。可分离落实为这条时间表。

能耗必须和中心熟化度误差 \(e_{\alpha}\) 一起读。单独看能耗，开环会因为提前停机、冷点还没起来，被误当成节能。

\section{蒸汽的三条路径与安全投影}
\label{ch:safe}

同一次蒸汽指令可以同时加热冷点、打湿表面，并在更长的时间上改变结壳热阻。把这三件事合成一个“蒸汽有利”的分数，会把薄件上的冷凝收益当作厚件上的菜单。决策状态的分量满足结构方程
\begin{equation}
\label{eq:scm}
Z_i=f_i(\mathrm{Pa}_i,U_i),\qquad i=1,\ldots,n_z.
\end{equation}
父节点 \(\mathrm{Pa}_i\) 由机理指定，见表\ref{tab:scm}。\(U_i\) 是外生噪声。学习阶段只调整函数的参数和事先声明的可选边；一次烹饪的短序列上，能量守恒的边保持固定。

\begin{table*}[t]
\centering
\caption{因果骨架。内部热阻既是参数也是子节点，因为结壳会增大它}
\label{tab:scm}
\setlength{\tabcolsep}{3.4pt}
{\small
\begin{tabularx}{\textwidth}{llY}
\toprule
子节点 & 父节点 & 物理内容 \\
\midrule
界面蒸汽压 & 蒸汽流、排汽、腔温 & 腔体水分平衡 \\
表面温度 & 蒸汽压、腔温、风速、加热、内部温度 & \eqref{eq:twonode} 第一式 \\
内部温度 & 表面温度、内部热阻、内部热容 & \eqref{eq:twonode} 第二式 \\
表面含水 & 蒸汽压、表面温度、传质系数 & 蒸发与冷凝 \\
冷点进度 & 内部温度的历史 & \eqref{eq:F} \\
表面褐变 & 表面温度、表面含水 & \eqref{eq:kinetics} \\
结壳 & 蛋白进度、表面含水 & \eqref{eq:structure} \\
内部热阻 & 结壳、厚度 & 结构到传递 \\
\bottomrule
\end{tabularx}}
\end{table*}

观测条件分布 \(p(\alpha_m\mid \dot m_{\mathrm{st}})\) 混有控制器对当前低温的反应：没熟所以多喷蒸汽。干预分布 \(p(\alpha_m\mid\mathrm{do}(\dot m_{\mathrm{st}}=m^\star))\) 把决定蒸汽流量的策略从结构方程中删去，只保留蒸汽到界面蒸汽压的物理边\cite{pearl2009}。孪生的积分实现干预：替换执行器轨迹后再积分。

蒸汽对品质的全导数分成冷点项、表面项和慢热阻项：
\begin{align}
\text{A: }& \dot m_{\mathrm{st}}\to q_{\mathrm{lat}}\to T_2\to\alpha_m\to Q,\\
\text{B: }& \dot m_{\mathrm{st}}\to\varepsilon_{l,s}\to(\alpha_b,m_{\mathrm{loss}})\to Q,\\
\text{C: }& \dot m_{\mathrm{st}}\to S\to R_{\mathrm{int}}\to T_2\to Q.
\end{align}
路径 A 在冷点落后时为正，路径 B 在表面已经湿且后续需要烤色时为负；路径 C 是慢路径，短探针窗内增量是 \(O(\tau/\tau_S)\)，全时域必须保留。总效应
\begin{equation}
\label{eq:path}
\begin{aligned}
\frac{\mathrm{d}Q}{\mathrm{d}m_{\mathrm{st}}}
={}&\underbrace{\frac{\partial Q}{\partial\alpha_m}\frac{\partial\alpha_m}{\partial m_{\mathrm{st}}}\Big|_{\!S}}_{\text{A}}\\
&+\underbrace{\frac{\partial Q}{\partial\varepsilon_{l,s}}\frac{\partial\varepsilon_{l,s}}{\partial m_{\mathrm{st}}}}_{\text{B}}\\
&+\underbrace{\frac{\partial Q}{\partial\alpha_m}\frac{\partial\alpha_m}{\partial R_{\mathrm{int}}}\frac{\partial R_{\mathrm{int}}}{\partial m_{\mathrm{st}}}}_{\text{C}}.
\end{aligned}
\end{equation}

\begin{proposition}[路径加性]
\label{prop:pathadd}
设品质泛函 \(Q\) 只通过 \(\alpha_m\)、\(\varepsilon_{l,s}\) 与 \(R_{\mathrm{int}}\) 三个量依赖蒸汽质量，且三者中只有
\(\alpha_m\) 与 \(R_{\mathrm{int}}\) 之间存在耦合。则全导数可写成\eqref{eq:path} 的三项，
其中第一项的 \(\partial\alpha_m/\partial m_{\mathrm{st}}\) 必须在固定 \(R_{\mathrm{int}}\) 下取值。
\end{proposition}

\begin{proof}
把 \(Q=Q(\alpha_m,\varepsilon_{l,s},R_{\mathrm{int}})\) 作为三个显式自变量的复合函数求全导数：
\[
\frac{\mathrm{d}Q}{\mathrm{d}m}
=\frac{\partial Q}{\partial\alpha_m}\frac{\mathrm{d}\alpha_m}{\mathrm{d}m}
+\frac{\partial Q}{\partial\varepsilon_{l,s}}\frac{\mathrm{d}\varepsilon_{l,s}}{\mathrm{d}m}
+\frac{\partial Q}{\partial R_{\mathrm{int}}}\frac{\mathrm{d}R_{\mathrm{int}}}{\mathrm{d}m}.
\]
其中 \(\alpha_m\) 本身又经 \(R_{\mathrm{int}}\to T_2\to\alpha_m\) 依赖蒸汽，故
\[
\frac{\mathrm{d}\alpha_m}{\mathrm{d}m}
=\frac{\partial\alpha_m}{\partial m}\Big|_{R_{\mathrm{int}}}
+\frac{\partial\alpha_m}{\partial R_{\mathrm{int}}}\frac{\mathrm{d}R_{\mathrm{int}}}{\mathrm{d}m}.
\]
代回并合项，即得\eqref{eq:path}。若第一项不在固定 \(R_{\mathrm{int}}\) 下取值，
则 \(\alpha_m\) 经 \(R_{\mathrm{int}}\) 的那条通道会被计入两次。
\end{proof}

这处固定条件是命题成立的关键。用一个具体的品质泛函核对：取
\(\alpha_m\) 为 \(T_2\) 的 Logistic、\(\varepsilon_{l,s}=0.5m_{\mathrm{st}}\)、
\(R_{\mathrm{int}}=3(1+1.2m_{\mathrm{st}})\)、\(Q=\alpha_m-0.5\varepsilon_{l,s}+0.1R_{\mathrm{int}}\)，
在 \(m_{\mathrm{st}}=1\) 处数值全导数为 \(+0.110000\)，三项之和同为 \(+0.110000\)，
残差 \(1.1\times10^{-10}\)。这条核对写在 \texttt{scripts/robust\_proxy.py} 的同源检验里；
命题\ref{prop:pathadd} 因此不是形式上的重新排列，而是一条可以逐项对账的分解。

\begin{figure}[htbp]
\centering
\colimg{S_paths.png}
\caption{同一次蒸汽指令的三条路径在核里的实际通量。路径 A 与 B 量级相当且符号相反，路径 C 小三个量级（图上放大一千倍），这是路径 C 在评分核里近乎不闭合的直接证据}
\label{fig:path}
\end{figure}

\begin{figure}[htbp]
\centering
\colimg{M_paths.png}
\caption{蒸汽离开界面之后的三条去向。路径 C 在核里不乘回导热}
\label{fig:mpaths}
\end{figure}

图\ref{fig:path}给出分叉，图\ref{fig:mpaths}把同一分叉画成三条去向。积分闭合了潜热通量对蒸汽、活度和风机的依赖，以及表面失水对活度和风机的依赖。结壳只累加，不乘进导热，所以路径 C 在评分核里没有闭合。界面通量对蒸汽指令的偏导数
\begin{equation}
\label{eq:dqlat}
\begin{aligned}
\frac{\partial q_{\mathrm{lat}}}{\partial u_{\mathrm{st}}}
={}&18000\,(1-a_w)\\
&+9000\max(a_w-0.35,0)\,u_v
\end{aligned}
\end{equation}
在饱和活度附近为正：关掉蒸汽时蒸发冷却被撤掉，蒸汽的热效应是撤冷却。短窗内路径 A 与 C 的温度增量很小。厚度增大使 A 变小，结壳使 B 更负；解释输出应当是这三项的符号与区间，而不是一个“蒸汽有利”的标量。

品质层的输出记为 \(\vect{u}_{\mathrm{AI}}\)。安全层求解
\begin{equation}
\label{eq:cbf}
\begin{aligned}
\vect{u}={}&\argmin_{\vect{u}}\;\|\vect{u}-\vect{u}_{\mathrm{AI}}\|_W^2\\
\mathrm{s.t.}\quad
&\dot h_j+\gamma_j h_j\ge 0.
\end{aligned}
\end{equation}
\(h_j\ge 0\) 表示安全集内部\cite{ames2017cbf}。自变量只含漏电、加热器温度、水位或空烧温升、腔压、门状态和执行器电流，不含 \(\alpha_m\)。屏障在估计器崩溃时仍然成立。漏电与干烧另有独立断电回路，不读取品质后验，也不经过产生参考动作的处理器。投影负责提前限制，断电负责投影失效后的切断。\cite{iec60335,gb4706}。

本次落实的削减是\eqref{eq:hard}：加热器温度超过 \(230^\circ\mathrm{C}\) 时把加热请求剪到 \(0.25\)。可行域与后验不相交——同一削减在两组不同后验上给出同一指令。置信度低于阈值时，送入投影的参考改为预先存储的低功率保持 \(0.15\)，学习模块的最后一次输出不再进入执行器。品质回路读表面温度和平均含水，写出四个请求；安全回路读水位、门、剩余电流和加热器温度，在请求与指令之间改写这四个量。投影可以减小加热，断电可以把四个量同时置零。图\ref{fig:d19proj}把本次工作点上的请求列与指令列画在一起：满功率平衡点 \(196.5^\circ\mathrm{C}\) 低于 \(230^\circ\mathrm{C}\)，两条曲线因此重合，投影处在待命状态——这正是“规则与后验无关”的可视形式，它不动作不是因为估计得好，而是因为这一次边界没有碰到它。

\begin{figure}[htbp]
\centering
\colimg{S_projection.png}
\caption{加热请求与指令。两条曲线全程重合，安全投影在本次工作点上待命；原因是满功率平衡点 $196.5^\circ\mathrm{C}$ 低于 $230^\circ\mathrm{C}$ 阈值}
\label{fig:d19proj}
\end{figure}

\section{模型假设与工程简化清单}

数值核沿厚度取十六条控制体，时间步 \(2\,\mathrm{s}\)。腔体压力场 \(p_g\) 与蒸汽质量分数 \(\omega_v\) 不作为独立未知量，界面潜热由蒸汽指令与活度的工程式\eqref{eq:qlat}给出。跟踪器为标量一阶滞后，时间常数 \(\tau=\mathrm{clip}(\hat r,0.8,8)\,(L/0.03)^2\cdot 280\,\mathrm{s}\)，先验 \(\hat r=5\)、方差 \(0.64\)。探针窗固定为 \([150,190)\,\mathrm{s}\)：前 \(20\,\mathrm{s}\) 蒸汽全开，后 \(20\,\mathrm{s}\) 风机升至 \(1.3\)；窗结束时若新热阻更接近由真值温差构成的参照 \(r_{\mathrm{proxy}}\)，方差乘 \(0.35\)。停机门槛为估计冷点越过 \(75^\circ\mathrm{C}\) 加随方差与厚度放大的裕量。加热器温度屏障只读管温，阈值 \(230^\circ\mathrm{C}\)；满功率平衡点约 \(196.5^\circ\mathrm{C}\)，主对照格子上削减保持待命。结构方程的路径积分与有限时域模型预测控制不进入评分循环；评分读取的是标量跟踪器给出的停机偏差。两节点教师模型的均方根误差约 \(41^\circ\mathrm{C}\)，留在规格层。褐变通道要求表面温度高于 \(110^\circ\mathrm{C}\) 且表面含水低于 \(0.55\)；主对照格子初始含水位于 \(0.64\)--\(0.80\)，表面保持湿润，褐变进度接近零。路径 C 经结壳累加，核未把结壳乘回导热。

下一章把上述停机几何放到同一积分器上对照：主格子、档位研究、网格加密、因子消融、宽格子与典型轨迹共用第\ref{ch:physics}章的工程常数和第\ref{ch:dual}章的跟踪器、探针窗与屏障。
\FloatBarrier


\section{四条策略在轨迹上的差别}
\label{sec:trace}

停机几何定的是"什么时候停"，轨迹形态定的是"停下来之前发生了什么"。两件事必须分开看：
一个策略可以在正确的时刻停下，却在过程中把表面烤干；也可以在过程中温和，却停得太早。
图\ref{fig:trpol}把四条策略放在同一初值上，画的是冷点、进度与加热这三条曲线。

\begin{figure*}[t]
\centering
\spanimg{T_policy.png}
\caption{种子 7、\(30\,\mathrm{mm}\)、组合档：四条策略的冷点、进度与加热。只有对偶把方差降到 \(0.224\)}
\label{fig:trpol}
\end{figure*}

\begin{figure}[htbp]
\centering
\colimg{T_thick.png}
\caption{种子 7、对偶、组合档：\(15\)、\(30\)、\(50\,\mathrm{mm}\) 的进度。\(50\,\mathrm{mm}\) 在 \(55\,\mathrm{min}\) 内停在 \(0.068\)}
\label{fig:trthick}
\end{figure}

\begin{figure*}[t]
\centering
\spanimg{T_cook.png}
\caption{种子 7、对偶、\(30\,\mathrm{mm}\)：蒸、烤、组合。蒸档进度停在 \(0.021\)，方差留在 \(0.64\)}
\label{fig:trcook}
\end{figure*}

图\ref{fig:trpol}里最值得注意的不是终点，而是中途。开环在 \(1170\,\mathrm{s}\) 被菜单切断，
冷点刚升到 \(16^\circ\mathrm{C}\) 就停止加热，其后曲线回落——停机偏差 \(e_{\alpha}=-0.849\)
并不是"算错了时刻"，而是这个时刻根本没有和内部状态发生关系。腔温比例则相反：它把加热一路
开到 \(3300\,\mathrm{s}\)，进度早已越过设计线，方差却始终停在先验 \(0.64\)。对偶在
\(192\,\mathrm{s}\) 把方差降到 \(0.224\)，\(1866\,\mathrm{s}\) 停机，三条曲线在停机附近同时收敛，
这是"信息动作改变了后验、后验改变了停机"唯一能在一张图上看出来的地方。

图\ref{fig:trthick}固定对偶与组合档，把厚度从 \(15\) 拉到 \(50\,\mathrm{mm}\)。薄件在时域内把
进度推到 \(1\)，停机由保持逻辑接管；厚件在 \(3300\,\mathrm{s}\) 时进度只有 \(0.068\)，停机时刻
被时域终点截住。三条曲线的结束方差都是 \(0.224\)：探针在三种厚度上都接受了新热阻，厚件的问题
不在辨识，而在扩散时间。

图\ref{fig:trcook}对比蒸、烤、组合三档。蒸档把加热夹在 \(0.25\) 以下，进度从一开始就走不动，
对偶在蒸档的中位偏差是 \(-0.830\)；烤档关掉蒸汽后只能靠加热器，表面虽然干爽，冷点却起得慢；
组合档是三者中唯一能同时给冷点潜热与表面排湿的允许集合。三档的差别不是策略的差别，
而是"允许集合"的差别，这一点在结果章会以交叉表的形式再次出现。
